\anexo{Modelos de dados utilizados na etapa de estruturação}
\label{an:modelos-pydantic}

Este anexo apresenta os modelos de dados desenvolvidos com a biblioteca \textit{Pydantic} e utilizados na etapa de estruturação das informações extraídas das transcrições de reuniões.

Os modelos definem formalmente os esquemas de saída empregados pelos modelos de linguagem por meio da funcionalidade de \textit{structured output}, garantindo padronização, validação automática e consistência semântica dos dados estruturados referentes a participantes, pautas, discussões e deliberações.

Cada modelo especifica explicitamente os campos esperados, seus tipos, descrições semânticas e exemplos ilustrativos, servindo tanto como contrato de saída quanto como mecanismo de validação dos resultados gerados.
\\
\\
\\

\subsection*{Modelo de Participantes}

\begin{lstlisting}[caption={Modelo Pydantic para representação de participantes}]
class Pessoa(BaseModel):
    """Representa uma única pessoa com nome e função."""
    nome: str = Field(..., description="O nome completo do participante.")
    cargo: str = Field(
        default="Não identificado",
        description="O cargo ou função do participante."
    )

class ParticipantesEncontradosDetalhado(BaseModel):
    participantes: List[Pessoa]
\end{lstlisting}

\pagebreak

\subsection*{Modelo de Pautas}

\begin{lstlisting}[caption={Modelo Pydantic para estruturação das pautas}]
class Pautas(BaseModel):
    titulo: str
    descricao: str
    responsavel: str = ""

class PautasEncontradas(BaseModel):
    pautas: List[Pautas]
\end{lstlisting}


\subsection*{Modelo de Discussões}

\begin{lstlisting}[caption={Modelo Pydantic para estruturação das discussões}]
class Discussao(BaseModel):
    assunto: str
    contexto: str

class DiscursoesEncontradas(BaseModel):
    discussoes: List[Discussao]
\end{lstlisting}


\subsection*{Modelo de Deliberações}

\begin{lstlisting}[caption={Modelo Pydantic para estruturação das deliberações}]
class Deliberacao(BaseModel):
    assunto: str
    acao: str
    responsavel: str = "Não encontrado"
    prazo: str = "Não identificado"

class DeliberacoesEncontradas(BaseModel):
    deliberacoes: List[Deliberacao]
\end{lstlisting}
